% !TEX root = ../main.tex
\section{Ejercicios resueltos}
El archivo \texttt{python/ejercicios\_resueltos.py} contiene los tres ejercicios ejecutables y comprueba los valores con aserciones.

\subsection{Resuelto 1: división de prueba y Fermat}
La división de prueba trata $n<2$, $2$, $3$, pares y divisores impares hasta \texttt{isqrt(n)}. Fermat verifica primero $\gcd(a,n)$ y después computa $a^{n-1}\bmod n$. El siguiente código utiliza las funciones de \texttt{primality.py} y exhibe los tres casos pedidos:
\begin{lstlisting}[language=Python,caption={División de prueba y Fermat: código ejecutable}]
def resuelto_1():
    """División de prueba y contraejemplo al recíproco de Fermat."""
    division_221 = trial_division(221)
    fermat_15 = fermat_round(15, 2)
    fermat_341 = fermat_round(341, 2)
    assert division_221 is False and fermat_15 is False and fermat_341 is True
    return {"221_es_primo": division_221, "15_pasa_fermat": fermat_15,
            "341_pasa_fermat": fermat_341}
\end{lstlisting}
\noindent\textbf{Resultado de la ejecución:}
\begin{lstlisting}[style=consolestyle]
>>> resuelto_1()
{'221_es_primo': False, '15_pasa_fermat': False, '341_pasa_fermat': True}
\end{lstlisting}

Aquí $341=11\cdot31$ pasa Fermat base $2$, aunque es compuesto. Para $561$, la base $2$ también pasa Fermat por ser Carmichael; la base $3$ comparte factor con $561$ y certifica composición.

\subsection{Resuelto 2: Miller--Rabin con traza verificable}
El par $(s,d)$ de $221-1=220=2^2\cdot55$ y la cadena de residuos de la base $2$ son:
\begin{equation}
2^{55}\bmod221=128,\qquad128^2\bmod221=30\neq220.
\end{equation}
No aparece $-1\equiv220$: la base $2$ demuestra que $221$ es compuesto. La salida incluye $n$, $a$, $s$, $d$, residuos, motivo y contadores de operaciones:
\begin{lstlisting}[language=Python,caption={Miller--Rabin con traza: código ejecutable}]
def resuelto_2():
    """Traza verificable de Miller--Rabin para n=221, base a=2."""
    assert decompose(221) == (2, 55)
    passed, detail = miller_rabin_round(221, 2, trace=True)
    assert not passed and detail["residuos"] == [128, 30]
    assert (detail["n"], detail["a"], detail["s"], detail["d"]) == (221, 2, 2, 55)
    return {"pasa": passed, **detail}
\end{lstlisting}
\noindent\textbf{Resultado de la ejecución:}
\begin{lstlisting}[style=consolestyle]
>>> resuelto_2()
{'pasa': False,
 'n': 221,
 'a': 2,
 's': 2,
 'd': 55,
 'residuos': [128, 30],
 'cuadrados_cadena': 1,
 'cuadrados_modulares': 6,
 'productos_modulares': 5,
 'motivo': 'sin_menos_uno'}
\end{lstlisting}

Se obtuvieron además las trazas exigidas: para $(n,a)=(341,2)$, $s=2$, $d=85$, cadena $[32,1]$; para $(561,2)$, $s=4$, $d=35$, cadena $[263,166,67,1]$. Ambas se rechazan por una raíz no trivial de $1$. En $341$, $\gcd(32-1,341)=31$ y $\gcd(32+1,341)=11$. Todas las trazas, también las de $1105$, $1729$ y un mentiroso fuerte de $25$, están en \texttt{results/trazas\_miller\_rabin.csv}.

\subsection{Resuelto 3: prueba repetida y candidatos}
\begin{lstlisting}[language=Python,caption={Prueba repetida y candidatos: código ejecutable}]
def resuelto_3():
    """Candidato reproducible y cinco rondas sobre el primo conocido 1009."""
    rng = Random(2026)
    n = generate_candidate(32, rng)
    assert n.bit_length() == 32 and n % 2 == 1
    report = primality_report(1009, rounds=5, rng=Random(2026))
    assert report["resultado"] == "probable_primo"
    assert report["rondas_ejecutadas"] == 5 and report["cota_error"] == 4.0**-5
    return {"candidato": n, "bits": n.bit_length(), "impar": bool(n % 2),
            "1009_resultado": report["resultado"],
            "rondas_ejecutadas": report["rondas_ejecutadas"],
            "cota_error": report["cota_error"]}
\end{lstlisting}
\noindent\textbf{Resultado de la ejecución:}
\begin{lstlisting}[style=consolestyle]
>>> resuelto_3()
{'candidato': 2659099673,
 'bits': 32,
 'impar': True,
 '1009_resultado': 'probable_primo',
 'rondas_ejecutadas': 5,
 'cota_error': 0.0009765625}
\end{lstlisting}

Los bits superior e inferior se fijan a $1$, respectivamente para garantizar longitud exacta e imparidad. Se eligió el primo conocido $1009$, mayor que los primos del precribado, para que se ejecuten realmente las cinco rondas: su resultado es ``probable primo'' y la cota condicional es $4^{-5}=1/1024$. Con semilla fija la ejecución es reproducible, pero \emph{no} segura para generar claves.
